% ch10 第10章习题解答（10 题）
\chapter{弦凝聚——光与费米子的统一}

\section*{问题 10.2.1}
\begin{problem}
（a）证明式 (10.2.1) 中的诸 $\hat U_{ij}$ 构成一组彼此对易的算符。（b）证明式 (10.2.3) 中的诸 $\psi$ 满足复费米子的标准反对易关系。
\end{problem}
\solution
\textbf{（a）预备与约定。}每格点四个马约拉纳算符满足
\begin{equation*}
\{\lambda_{\pmb{i}}^{a},\lambda_{\pmb{j}}^{b}\}=2\delta_{ab}\delta_{\pmb{i}\pmb{j}},
\qquad a,b=x,\bar{x},y,\bar{y},
\end{equation*}
即：同格点不同味的马约拉纳算符反对易，不同格点的对易，且 $(\lambda_{\pmb{i}}^{a})^2=1$。链算符按正文取（我们把指向 $+\hat{x},+\hat{y}$ 的链写作
\begin{equation*}
\hat U_{\pmb{i},\pmb{i}+\hat{\pmb{x}}}=\mathrm{i}\,\lambda_{\pmb{i}}^{x}\lambda_{\pmb{i}+\hat{\pmb{x}}}^{\bar{x}},\qquad
\hat U_{\pmb{i},\pmb{i}+\hat{\pmb{y}}}=\mathrm{i}\,\lambda_{\pmb{i}}^{y}\lambda_{\pmb{i}+\hat{\pmb{y}}}^{\bar{y}},\qquad
\hat U_{\pmb{j}\pmb{i}}\equiv-\hat U_{\pmb{i}\pmb{j}} .
\end{equation*}
\emph{（译注：原书式 (10.2.1) 印作 $\hat U_{\pmb{i},\pmb{i}+\hat{\pmb x}}=\lambda_{\pmb i}^{x}\lambda_{\pmb{i+\hat x}}^{\bar x}$，无因子 $\mathrm{i}$；按该定义 $(\hat U_{\pmb{ij}})^2=+1$，与正文"$\hat U_{\pmb{ij}}^2=-1$、本征值 $s_{\pmb{ij}}=\pm\mathrm{i}$"矛盾。按本意应如上补一因子 $\mathrm{i}$（或在谱值约定中等价地取 $s_{\pmb{ij}}=\pm\mathrm{i}$）。该相因子不影响任何（反）对易关系，下面全部结论与此约定无关。）}

先证两条基本性质：
\begin{enumerate}
\item[(i)] $\hat U_{\pmb{ij}}^2=-1$。事实上
\begin{equation*}
\hat U_{\pmb{ij}}^2=(\mathrm{i}\lambda_{\pmb{i}}^{a}\lambda_{\pmb{j}}^{b})(\mathrm{i}\lambda_{\pmb{i}}^{a}\lambda_{\pmb{j}}^{b})
=-\,\lambda_{\pmb{i}}^{a}\lambda_{\pmb{j}}^{b}\lambda_{\pmb{i}}^{a}\lambda_{\pmb{j}}^{b}
=-\,\lambda_{\pmb{i}}^{a}\lambda_{\pmb{i}}^{a}\lambda_{\pmb{j}}^{b}\lambda_{\pmb{j}}^{b}=-1,
\end{equation*}
其中把第二个 $\lambda_{\pmb{i}}^{a}$ 移过 $\lambda_{\pmb{j}}^{b}$（不同格点，对易）不产生符号。故 $\hat U_{\pmb{ij}}$ 的本征值为 $s_{\pmb{ij}}=\pm\mathrm{i}$，与 $s_{\pmb{ij}}=-s_{\pmb{ji}}$（由 $\hat U_{\pmb{ji}}=-\hat U_{\pmb{ij}}$）一致。
\item[(ii)] 不共享格点的两条链算符对易：设 $\{\pmb{i},\pmb{j}\}\cap\{\pmb{k},\pmb{l}\}=\varnothing$，则 $\hat U_{\pmb{ij}}$ 的两个马约拉纳算符与 $\hat U_{\pmb{kl}}$ 的两个全部不同格点对易，直接换序无符号。
\end{enumerate}

\textbf{共享格点的两条链。}设两条链都从 $\pmb{i}$ 出发：$\hat U_1=\mathrm{i}\lambda_{\pmb{i}}^{a}\lambda_{\pmb{j}}^{b}$，$\hat U_2=\mathrm{i}\lambda_{\pmb{i}}^{c}\lambda_{\pmb{k}}^{d}$（$\pmb{j}\neq\pmb{k}$）。把两个乘积都化到规范次序（格点 $\pmb i$ 的两个算符在前、按味排序，$\pmb j,\pmb k$ 的在后）：
\begin{align*}
\hat U_1\hat U_2&=-\,\lambda_{\pmb{i}}^{a}\lambda_{\pmb{j}}^{b}\lambda_{\pmb{i}}^{c}\lambda_{\pmb{k}}^{d}
=-\,\lambda_{\pmb{i}}^{a}\lambda_{\pmb{i}}^{c}\lambda_{\pmb{j}}^{b}\lambda_{\pmb{k}}^{d},\\
\hat U_2\hat U_1&=-\,\lambda_{\pmb{i}}^{c}\lambda_{\pmb{k}}^{d}\lambda_{\pmb{i}}^{a}\lambda_{\pmb{j}}^{b}
=-\,\lambda_{\pmb{i}}^{c}\lambda_{\pmb{i}}^{a}\lambda_{\pmb{j}}^{b}\lambda_{\pmb{k}}^{d}
=+\,\lambda_{\pmb{i}}^{a}\lambda_{\pmb{i}}^{c}\lambda_{\pmb{j}}^{b}\lambda_{\pmb{k}}^{d},
\end{align*}
第一步只用了不同格点算符的对易性，第二步用了 $\lambda_{\pmb{i}}^{c}\lambda_{\pmb{i}}^{a}=-\lambda_{\pmb{i}}^{a}\lambda_{\pmb{i}}^{c}$（$a\neq c$）。于是
\begin{equation*}
\hat U_1\hat U_2=-\,\hat U_2\hat U_1\qquad(a\neq c);
\end{equation*}
而若 $a=c$，则两算符共享同一个马约拉纳算符 $\lambda_{\pmb{i}}^{a}$：
\begin{equation*}
\hat U_1\hat U_2=(\mathrm{i})^2\lambda_{\pmb{i}}^{a}\lambda_{\pmb{j}}^{b}\lambda_{\pmb{i}}^{a}\lambda_{\pmb{k}}^{d}
=\lambda_{\pmb{j}}^{b}\lambda_{\pmb{k}}^{d}=\hat U_2\hat U_1 ,
\end{equation*}
即对易。

对式 (10.2.1) 的味分配，格点 $\pmb{i}$ 的四条链（出 $\hat x$、入 $\hat x$、出 $\hat y$、入 $\hat y$）分别携带 $\lambda_{\pmb i}^{x},\lambda_{\pmb i}^{\bar x},\lambda_{\pmb i}^{y},\lambda_{\pmb i}^{\bar y}$——同一格点的两条不同链携带不同味，故
\begin{equation*}
\{\hat U_{\pmb{ij}},\hat U_{\pmb{ik}}\}=0\qquad(\text{相邻链}),\qquad
[\hat U_{\pmb{ij}},\hat U_{\pmb{kl}}]=0\qquad(\text{不相邻链}).
\end{equation*}
\emph{（译注：因此，严格地说诸 $\hat U_{\pmb{ij}}$ 并不全彼此对易——原书（及 PRL 90, 016803）"诸 $\hat U_{\pmb{ij}}$ 构成对易集合"的断言欠准确。可解性真正需要、并且确实成立的是下面三条：(i) 方格通量算符 $\hat F_{\pmb i}$ 彼此对易；(ii) 每个 $\hat F_{\pmb i}$ 与每个 $\hat U_{\pmb{ij}}$ 对易（因而 $H=-\sum V_{\pmb i}\hat F_{\pmb i}$ 与所有 $\hat U_{\pmb{ij}}$ 对易）；(iii) $\hat F_{\pmb i}^2=1$。证得这三条后，$H$ 由对易算符 $\{\hat F_{\pmb i}\}$ 的共同本征态 $|\{F_{\pmb i}\}\rangle$ 对角化，即得精确解（10.2.4）；用链变量组态 $\{s_{\pmb{ij}}\}$ 标记态时，则须如正文 10.2.2 节那样按 $Z_2$ 规范等价类与奇偶约束理解。）}

\textbf{通量算符的对易性。}用上面的计算直接给出方格算符。注意 $\hat U_{\pmb i_2\pmb i_3}=-\mathrm{i}\lambda_{\pmb i_3}^{x}\lambda_{\pmb i_2}^{\bar x}$（$\pmb i_3\to\pmb i_2$ 为 $+x$ 方向）、$\hat U_{\pmb i_3\pmb i}=-\mathrm{i}\lambda_{\pmb i}^{y}\lambda_{\pmb i_3}^{\bar y}$，于是
\begin{align*}
\hat F_{\pmb i}&=\hat U_{\pmb i\pmb i_1}\hat U_{\pmb i_1\pmb i_2}\hat U_{\pmb i_2\pmb i_3}\hat U_{\pmb i_3\pmb i}\\
&=-\big(\lambda_{\pmb i}^{x}\lambda_{\pmb i}^{y}\big)\big(\lambda_{\pmb i_3}^{x}\lambda_{\pmb i_3}^{\bar y}\big)\big(\lambda_{\pmb i_1}^{\bar x}\lambda_{\pmb i_1}^{y}\big)\big(\lambda_{\pmb i_2}^{\bar x}\lambda_{\pmb i_2}^{\bar y}\big),
\end{align*}
即 $\hat F_{\pmb i}$ 是四个"角因子"（每个角上一对马约拉纳算符，分别来自该角的两个方向链）之积再乘 $(-1)$（相乘次序不同只差整体符号，不影响下述论证）。
\begin{enumerate}
\item[(iii)] $\hat F_{\pmb i}^2=1$：把 $\hat F_{\pmb i}$ 的八个马约拉纳算符逐一平方并换序，同格点对换共 $4$ 次（每个角因子内部一对），符号 $(-1)^4=+1$，得 $\hat F_{\pmb i}^2=+1$。
\item[(iv)] $[\hat F_{\pmb i},\hat F_{\pmb j}]=0$。关键计数：同一格点上两对马约拉纳算符（各自为"双线性"因子）换序时，仅当两对共享一个马约拉纳算符才产生 $(-1)$；不共享时换序需 $2\times2=4$ 次同格点对换，符号 $+1$；完全相同则平方为 $1$。两个方格 $\hat F_{\pmb i},\hat F_{\pmb j}$ 要么不相交（平凡对易），要么共享两个格点：
\begin{itemize}
\item 共享一条边（如 $\pmb j=\pmb i_1$）：在共享的两个格点上，两方格的角因子各共享一个马约拉纳算符（公共边的味），每个共享格点贡献 $(-1)$，共 $(-1)^2=+1$；
\item 只共享一个对角格点（$\pmb j=\pmb i_2$）：该格点上两角因子的四个马约拉纳算符来自四条不同的链，互不相同，换序符号 $+1$。
\end{itemize}
故恒有 $[\hat F_{\pmb i},\hat F_{\pmb j}]=0$。同理，链算符 $\hat U_{\pmb{k}\pmb l}$ 与 $\hat F_{\pmb i}$ 至多共享两个格点、每处恰共享一个马约拉纳算符，故 $[\hat F_{\pmb i},\hat U_{\pmb{k}\pmb l}]=0$，从而 $[H,\hat U_{\pmb{k}\pmb l}]=0$。
\end{enumerate}
综上，模型由对易的通量算符系统精确可解：$H=-\sum_{\pmb i}V_{\pmb i}\hat F_{\pmb i}$ 的本征态为 $|\{F_{\pmb i}\}\rangle$，能量 $E=-\sum_{\pmb i}V_{\pmb i}F_{\pmb i}$，$F_{\pmb i}=\pm1$。

\textbf{（b）}由 $2\psi_{1}=\lambda^{x}+\mathrm{i}\lambda^{\bar x}$ 得 $2\psi_{1}^{\dagger}=\lambda^{x}-\mathrm{i}\lambda^{\bar x}$。逐项计算：
\begin{equation*}
\{\psi_1,\psi_1^{\dagger}\}=\frac14\{\lambda^{x}+\mathrm{i}\lambda^{\bar x},\lambda^{x}-\mathrm{i}\lambda^{\bar x}\}
=\frac14\big[\{\lambda^{x},\lambda^{x}\}+\{\lambda^{\bar x},\lambda^{\bar x}\}\big]
=\frac14(2+2)=1,
\end{equation*}
（交叉项 $\{\lambda^{x},\lambda^{\bar x}\}=0$）。同理 $\{\psi_2,\psi_2^{\dagger}\}=1$。对交叉项，$\psi_1$ 只含 $\lambda^{x},\lambda^{\bar x}$，$\psi_2$ 只含 $\lambda^{y},\lambda^{\bar y}$，四对交叉反对易子全为零，例如
\begin{equation*}
\{\psi_1,\psi_2^{\dagger}\}=\frac14\{\lambda^{x}+\mathrm{i}\lambda^{\bar x},\lambda^{y}-\mathrm{i}\lambda^{\bar y}\}=0,
\end{equation*}
（同格点不同味反对易给出 $\{x,y\}=\{x,\bar y\}=\{\bar x,y\}=\{\bar x,\bar y\}=0$）。故
\begin{equation*}
\boxed{\;\{\psi_{a,\pmb i},\psi_{b,\pmb j}^{\dagger}\}=\delta_{ab}\delta_{\pmb{i}\pmb j},\qquad
\{\psi_{a,\pmb i},\psi_{b,\pmb j}\}=0\;}
\end{equation*}
即 $\psi_{1,2}$ 是标准的复费米子。每个格点上由它们生成 $4$ 维希尔伯特空间（费米子数 $0,1,2$），全空间维数 $4^{N_{\text{site}}}$。

\section*{问题 10.2.2}
\begin{problem}
（a）证明 $2\psi_1^\dagger\psi_1-1=\mathrm{i}\lambda^{x}\lambda^{\bar{x}}$、$2\psi_2^\dagger\psi_2-1=\mathrm{i}\lambda^{y}\lambda^{\bar{y}}$，并证明一个格点上具有偶数个费米子的态是 $\lambda^{x}\lambda^{y}\lambda^{\bar{x}}\lambda^{\bar{y}}$ 的本征值为 $+1$ 的本征态。（b）证明式 (10.2.7)。（c）证明式 (10.2.5) 中的诸 $\sigma^{l}$ 在物理希尔伯特空间内满足泡利矩阵代数。（d）证明费米子哈密顿量 (10.2.2) 在物理希尔伯特空间内变为自旋哈密顿量 (10.2.6)。
\end{problem}
\solution
\textbf{（a）}由 $2\psi_1=\lambda^{x}+\mathrm{i}\lambda^{\bar{x}}$（同格点，下面略去格点指标）：
\begin{equation*}
\psi_1^\dagger\psi_1=\frac14(\lambda^{x}-\mathrm{i}\lambda^{\bar{x}})(\lambda^{x}+\mathrm{i}\lambda^{\bar{x}})
=\frac14\big[1+\mathrm{i}\lambda^{x}\lambda^{\bar{x}}-\mathrm{i}\lambda^{\bar{x}}\lambda^{x}+1\big]
=\frac12\big(1+\mathrm{i}\lambda^{x}\lambda^{\bar{x}}\big),
\end{equation*}
（用了 $\lambda^{\bar{x}}\lambda^{x}=-\lambda^{x}\lambda^{\bar{x}}$、$(\lambda^{a})^2=1$）。故
\begin{equation*}
\boxed{\;2\psi_1^\dagger\psi_1-1=\mathrm{i}\lambda^{x}\lambda^{\bar{x}},\qquad 2\psi_2^\dagger\psi_2-1=\mathrm{i}\lambda^{y}\lambda^{\bar{y}}\;}
\end{equation*}
另一方面，粒子数奇偶算符为
\begin{equation*}
(-)^{n_1+n_2}=(2\psi_1^\dagger\psi_1-1)(2\psi_2^\dagger\psi_2-1)
=\big(\mathrm{i}\lambda^{x}\lambda^{\bar{x}}\big)\big(\mathrm{i}\lambda^{y}\lambda^{\bar{y}}\big)
=-\lambda^{x}\lambda^{\bar{x}}\lambda^{y}\lambda^{\bar{y}}
=\lambda^{x}\lambda^{y}\lambda^{\bar{x}}\lambda^{\bar{y}},
\end{equation*}
最后一步把 $\lambda^{\bar x}$ 移过 $\lambda^{y}$（同格点，$-1$ 号）。于是作为算符恒等式
\begin{equation*}
(-)^{\psi_1^\dagger\psi_1+\psi_2^\dagger\psi_2}=\lambda^{x}\lambda^{y}\lambda^{\bar{x}}\lambda^{\bar{y}},
\end{equation*}
偶费米子态（$n_1+n_2$ 偶，即 $(-)^{n_1+n_2}=+1$）正是 $\lambda^{x}\lambda^{y}\lambda^{\bar{x}}\lambda^{\bar{y}}$ 本征值为 $+1$ 的本征态。这就是物理希尔伯特空间的约束 $\lambda^{x}\lambda^{y}\lambda^{\bar{x}}\lambda^{\bar{y}}=1$。

\textbf{（b）}全部 $2N_{\text{site}}$ 条链算符的乘积中，每条链用掉其两端各一个马约拉纳算符；每个格点恰好有四条关联链（出 $\hat x$、入 $\hat x$、出 $\hat y$、入 $\hat y$），分别用掉 $\lambda_{\pmb i}^{x},\lambda_{\pmb i}^{\bar x},\lambda_{\pmb i}^{y},\lambda_{\pmb i}^{\bar y}$。因此乘积
\begin{equation*}
\prod_{\text{links}}\hat U_{\pmb{i}\pmb{j}}
\end{equation*}
中每个马约拉纳算符恰好出现一次。取各链算符相乘的次序使同格点换序的净符号为 $+1$（例如按格点字典序、每点先 $\hat x$ 链后 $\hat y$ 链并按规范次序排列；任何残余的 $\pm1$ 均可吸收进各链的相位约定，即原书隐含的约定），得算符恒等式
\begin{equation*}
\prod_{\text{links}}\hat U_{\pmb{i}\pmb{j}}=\prod_{\pmb i}\lambda_{\pmb i}^{x}\lambda_{\pmb i}^{y}\lambda_{\pmb i}^{\bar x}\lambda_{\pmb i}^{\bar y}=(-)^{\hat N_f},
\end{equation*}
第二步用了（a）中的恒等式。取 $|\{s_{\pmb{ij}}\}\rangle$ 的本征值（$\hat U_{\pmb{i}\pmb j}$ 的本征值为 $s_{\pmb{ij}}=\pm\mathrm{i}$，故每条链贡献 $\mathrm{i}s_{\pmb{i}\pmb j}=\mp1$），即得式 (10.2.7)：
\begin{equation*}
\boxed{\;\prod_{\pmb i}\big(\mathrm{i}s_{\pmb i,\pmb i+\pmb x}\big)\big(\mathrm{i}s_{\pmb i,\pmb i+\pmb y}\big)=(-)^{N_f}\;}
\end{equation*}
它把链变量组态的总"符号"与态的费米子数奇偶联系起来。对每格点偶数费米子的物理态，$(-)^{N_f}=+1$，立即得到投影非零条件 (10.2.8)：
\begin{equation*}
\prod_{\pmb i}\big(\mathrm{i}s_{\pmb i,\pmb i+\pmb x}\big)\big(\mathrm{i}s_{\pmb i,\pmb i+\pmb y}\big)=1 .
\end{equation*}

\textbf{（c）}定义 (10.2.5)：$\sigma^{x}=\mathrm{i}\lambda^{y}\lambda^{x}$，$\sigma^{y}=\mathrm{i}\lambda^{\bar x}\lambda^{y}$，$\sigma^{z}=\mathrm{i}\lambda^{x}\lambda^{\bar x}$（均含偶数个马约拉纳算符，故保持每格点费米子奇偶，作用在物理希尔伯特空间之内）。厄米性：如 $(\sigma^{x})^{\dagger}=-\mathrm{i}\,(\lambda^{x})^{\dagger}(\lambda^{y})^{\dagger}=-\mathrm{i}\lambda^{x}\lambda^{y}=\mathrm{i}\lambda^{y}\lambda^{x}=\sigma^{x}$，同理其余。代数：
\begin{align*}
\sigma^{x}\sigma^{y}&=(\mathrm{i})^{2}\lambda^{y}\lambda^{x}\lambda^{\bar{x}}\lambda^{y}
=-\lambda^{y}\lambda^{x}\lambda^{\bar{x}}\lambda^{y}
=-\lambda^{x}\lambda^{\bar{x}}
=\mathrm{i}\sigma^{z},\
\sigma^{y}\sigma^{z}&=-\lambda^{\bar{x}}\lambda^{y}\lambda^{x}\lambda^{\bar{x}}
=-\lambda^{y}\lambda^{x}
=\mathrm{i}\sigma^{x},\
\sigma^{z}\sigma^{x}&=-\lambda^{x}\lambda^{\bar{x}}\lambda^{y}\lambda^{x}
=-\lambda^{\bar{x}}\lambda^{y}
=\mathrm{i}\sigma^{y},
\end{align*}
（第一行：把最右 $\lambda^{y}$ 依次左移过 $\lambda^{\bar x},\lambda^{x}$ 两次同格点换序得 $\lambda^{x}\lambda^{\bar x}$；第二、三行类似）。由厄米性即得 $\{\sigma^{a},\sigma^{b}\}=2\delta_{ab}$，$(\sigma^{a})^2=1$。不同格点的 $\sigma$ 由不同格点的马约拉纳算符构成，彼此对易。这些代数关系是全希尔伯特空间上的算符恒等式；限制在二维的物理空间内，三个泡利矩阵满足完备代数、表示不可约。故诸 $\sigma^{l}$ 在物理希尔伯特空间内构成自旋 $1/2$ 算符的忠实表示。

\textbf{（d）}把四条链算符显式写出（用 $\hat U_{\pmb{i}\pmb j}=-\hat U_{\pmb{j}\pmb i}$：$\hat U_{\pmb i_2\pmb i_3}=-\mathrm{i}\lambda_{\pmb i_3}^{x}\lambda_{\pmb i_2}^{\bar x}$，$\hat U_{\pmb i_3\pmb i}=-\mathrm{i}\lambda_{\pmb i}^{y}\lambda_{\pmb i_3}^{\bar y}$；四个 $\mathrm i$ 相乘为 $+1$）：
\begin{equation*}
\hat F_{\pmb i}
=\hat U_{\pmb i\pmb i_1}\hat U_{\pmb i_1\pmb i_2}\hat U_{\pmb i_2\pmb i_3}\hat U_{\pmb i_3\pmb i}
=-\big(\lambda_{\pmb i}^{x}\lambda_{\pmb i}^{y}\big)\big(\lambda_{\pmb i_3}^{x}\lambda_{\pmb i_3}^{\bar y}\big)\big(\lambda_{\pmb i_1}^{\bar x}\lambda_{\pmb i_1}^{y}\big)\big(\lambda_{\pmb i_2}^{\bar x}\lambda_{\pmb i_2}^{\bar y}\big),
\end{equation*}
（规范次序下的四个"角双线性"因子）。在物理希尔伯特空间内利用（a）中的恒等式及其推论（由 $\lambda^{x}\lambda^{y}\lambda^{\bar x}\lambda^{\bar y}=1$ 移项：$\lambda^{\bar x}\lambda^{\bar y}=-\lambda^{x}\lambda^{y}$，$\lambda^{x}\lambda^{\bar y}=\lambda^{\bar x}\lambda^{y}$）：
\begin{equation*}
\lambda^{x}\lambda^{y}=\mathrm{i}\sigma^{x},\qquad
\lambda^{\bar{x}}\lambda^{y}=-\mathrm{i}\sigma^{y},\qquad
\lambda^{\bar{x}}\lambda^{\bar{y}}=-\mathrm{i}\sigma^{x},\qquad
\lambda^{x}\lambda^{\bar{y}}=-\mathrm{i}\sigma^{y}\qquad(\text{物理空间}),
\end{equation*}
代入（四个相因子的积 $=(\mathrm{i})(-\mathrm{i})(-\mathrm{i})(-\mathrm{i})=-1$）：
\begin{equation*}
\hat F_{\pmb i}\Big|_{\text{物理}}
=-(-1)\,\sigma_{\pmb i}^{x}\sigma_{\pmb i_3}^{y}\sigma_{\pmb i_1}^{y}\sigma_{\pmb i_2}^{x}
=\sigma_{\pmb i}^{x}\sigma_{\pmb i+\hat{\pmb x}}^{y}\sigma_{\pmb i+\hat{\pmb x}+\hat{\pmb y}}^{x}\sigma_{\pmb i+\hat{\pmb y}}^{y},
\end{equation*}
（不同格点的 $\sigma$ 对易，可任意排序）。故
\begin{equation*}
\boxed{\;H=-\sum_{\pmb i}V_{\pmb i}\hat F_{\pmb i}
=-\sum_{\pmb i}V_{\pmb i}\,\sigma_{\pmb i}^{x}\sigma_{\pmb i+\hat{\pmb x}}^{y}\sigma_{\pmb i+\hat{\pmb x}+\hat{\pmb y}}^{x}\sigma_{\pmb i+\hat{\pmb y}}^{y}
=H_{\text{spin}}\quad(\text{物理希尔伯特空间内})\;}
\end{equation*}
即式 (10.2.6)。注意 $\hat F_{\pmb i}$ 含偶数个马约拉纳算符于每个角，保持每格点费米子奇偶，因而 $H$ 分别作用在偶、奇两个子空间内。

\section*{问题 10.2.3}
\begin{problem}
把费米子哈密顿量 $H$ 限制在每格点具有奇数个费米子的子空间内，求相应的自旋 $1/2$ 系统。
\end{problem}
\solution
\textbf{奇子空间是不变子空间。}$\hat F_{\pmb i}$ 的表达式 (10.2.2) 表明它在每个方格的四个角上各含两个（偶数个）马约拉纳算符，故 $[(-)^{n_{\pmb i}},\hat F_{\pmb i}]=0$ 对所有 $\pmb i$ 成立，$H$ 保持"每格点奇数个费米子"的子空间不变。

\textbf{自旋算符。}$\sigma_{\pmb i}^{l}$（式 (10.2.5)）是偶算符，分别在偶、奇两个子空间内作用，且（c）小问中验证的泡利代数 $\sigma^{x}\sigma^{y}=\mathrm{i}\sigma^{z}$ 等是全空间上的算符恒等式，因而在奇子空间内同样成立；又奇子空间每格点二维，泡利代数在该二维空间上的表示不可约。故同一组 $\sigma^{l}$ 就是奇子空间的自旋算符。

\textbf{哈密顿量。}关键恒等式：在全希尔伯特空间上可以验证（把角因子规范化后比较）
\begin{equation*}
\sigma_{\pmb i}^{x}\sigma_{\pmb i_1}^{y}\sigma_{\pmb i_2}^{x}\sigma_{\pmb i_3}^{y}
=\Pi_{\pmb i_2}\Pi_{\pmb i_3}\,\hat F_{\pmb i},\qquad
\Pi_{\pmb i}\equiv\lambda_{\pmb i}^{x}\lambda_{\pmb i}^{y}\lambda_{\pmb i}^{\bar x}\lambda_{\pmb i}^{\bar y}=(-)^{n_{\pmb i}},
\end{equation*}
（其中 $\Pi_{\pmb i}=(-)^{n_{\pmb i}}$ 为（a）小问的奇偶算符；例如直接把 $\hat F_{\pmb i}$ 的四个角因子 $\lambda^{x}\lambda^{y}=\mathrm{i}\sigma^{x}$、$\lambda^{\bar x}\lambda^{y}=-\mathrm{i}\sigma^{y}$ 精确代入，而 $\lambda^{x}\lambda^{\bar y}=-\mathrm{i}\Pi\sigma^{y}$、$\lambda^{\bar x}\lambda^{\bar y}=-\mathrm{i}\Pi\sigma^{x}$ 含奇偶算符）。在偶子空间内 $\Pi_{\pmb i}=+1$，在奇子空间内每个格点 $\Pi_{\pmb i}=-1$；由于上式只含 $\Pi_{\pmb i_2}\Pi_{\pmb i_3}$ 两个奇偶算符的乘积，两个子空间中均给出
\begin{equation*}
\hat F_{\pmb i}\Big|_{\text{奇}}=\sigma_{\pmb i}^{x}\sigma_{\pmb i+\hat{\pmb x}}^{y}\sigma_{\pmb i+\hat{\pmb x}+\hat{\pmb y}}^{x}\sigma_{\pmb i+\hat{\pmb y}}^{y}\Big|_{\text{奇}} .
\end{equation*}
因此
\begin{equation*}
\boxed{\;H\big|_{\text{奇子空间}}=H_{\text{spin}}=-\sum_{\pmb i}V_{\pmb i}\,\sigma_{\pmb i}^{x}\sigma_{\pmb i+\hat{\pmb x}}^{y}\sigma_{\pmb i+\hat{\pmb x}+\hat{\pmb y}}^{x}\sigma_{\pmb i+\hat{\pmb y}}^{y}\;}
\end{equation*}
与偶子空间中的 (10.2.6) 形式完全相同：费米子系统的希尔伯特空间分解为偶、奇两个子空间，每个子空间都是同一个自旋 $1/2$ 方格模型（$H=H_{\text{spin}}\oplus H_{\text{spin}}$），能谱与简并结构相同（通量算符的全部代数关系，如 $\prod\hat F_{\pmb i}=1$ 等，只是 $\sigma$ 代数的恒等式，在两个表示中同样成立）。

两者的差别只在精确态的"链变量"标记上：由 (10.2.7)，偶子空间的态要求 $\prod_{\pmb i}(\mathrm{i}s_{\pmb i,\pmb i+\pmb x})(\mathrm{i}s_{\pmb i,\pmb i+\pmb y})=+1$，而奇子空间的态每格点奇数个费米子，$(-)^{N_f}=(-1)^{N_{\text{site}}}$，故约束改为
\begin{equation*}
\prod_{\pmb i}\big(\mathrm{i}s_{\pmb i,\pmb i+\pmb x}\big)\big(\mathrm{i}s_{\pmb i,\pmb i+\pmb y}\big)=(-1)^{N_{\text{site}}},
\end{equation*}
投影算符则换为 $\prod_{\pmb i}(1-\Pi_{\pmb i})/2$。在偶 $\times$ 偶格子上 $N_{\text{site}}$ 为偶，两个子空间由同一批组态 $\{s_{\pmb{ij}}\}$ 生成，是同一自旋模型的两份全同拷贝。

\section*{问题 10.2.4}
\begin{problem}
考虑 $L_x\times L_y$ 周期格子上的式 (10.2.6)，设 $V_{\pmb i}=V$，仅在某一行方格上 $V_{\pmb i}=0$。此时系统可视为定义在 $x$ 方向具有两条圆形边缘的圆柱上。求基态简并度；把这些简并基态视为边缘态，证明每个边缘格点有 $\sqrt{2}$ 个边缘态，即边缘激发由一个马约拉纳费米子描写。
\end{problem}
\solution
取 $V>0$（$V<0$ 时把下述论证中的 $F_{\pmb i}=+1$ 换成 $F_{\pmb i}=\mathrm{sgn}(V)=-1$，只要约束允许该通量组态，简并计数完全相同）。设被抽掉相互作用的一行方格为 $y=0$ 行（共 $L_x$ 个方格），并设格子为偶 $\times$ 偶（与 10.2.2 节一致）。

\textbf{基态条件。}基态要求所有 $V\neq0$ 的方格上 $F_{\pmb i}=+1$；缺口行方格的 $V_{\pmb i}=0$，其通量不受哈密顿量约束。由问题 10.3.1 中证明的恒等式 $\prod_{\text{全部}}\hat F_{\pmb i}=1$ 与 $\prod_{\text{偶方格}}\hat F_{\pmb i}=1$，缺口行通量须满足
\begin{equation*}
\prod_{x}F_{(x,0)}=1,\qquad \prod_{x\ \text{偶}}F_{(x,0)}=1 ,
\end{equation*}
两个独立约束，故缺口行自由的通量组态共 $2^{L_x-2}$ 个。

\textbf{简并度。}对每个固定通量组态，基态还携带圆柱拓扑简并：缺口行把环面分成两个半圆柱，每个半圆柱各有一条不可收缩的 $x$ 方向回路；仿照 10.2.2 节的 $T_1$ 操作（沿回路插入 $\pi$ 通量），在每个半圆柱上各得到一个二重简并，共 $4$ 重。于是
\begin{equation*}
\boxed{\;\text{基态简并度}=4\times2^{L_x-2}=2^{L_x}\;}
\end{equation*}
（与稳定子秩的独立计数一致）。这 $2^{L_x}$ 个基态可显式生成：跨越缺口行的 $L_x$ 条 $y$ 方向链变量，其翻转算符只与缺口行（$V=0$）的方格相邻，因而与全部 $V\neq0$ 的 $\hat F_{\pmb i}$ 对易；任选子集翻转即得不同基态，$L_x$ 个独立二元选择给出 $2^{L_x}$ 个态。

\textbf{每边缘格点 $\sqrt2$ 个态：马约拉纳诠释。}两条圆形边缘各有 $L_x$ 个格点。基态数按边缘分配
\begin{equation*}
2^{L_x}=2^{L_x/2}\times2^{L_x/2},
\end{equation*}
即每条边缘有 $2^{L_x/2}=(\sqrt2\,)^{L_x}$ 个态，
\begin{equation*}
\boxed{\;\text{每个边缘格点的态数}=2^{1/2}=\sqrt2\;}
\end{equation*}
这正是马约拉纳费米子零模的标志：若边缘的 $L_x$ 个零模为马约拉纳算符 $\gamma_1,\dots,\gamma_{L_x}$（$\{\gamma_m,\gamma_n\}=2\delta_{mn}$），则它们配对成 $L_x/2$ 个复费米子零模，占据数给出 $2^{L_x/2}$ 个态——每格点 $\sqrt2$ 个；反之，若每格点是一个复（狄拉克）零模，则应有 $2^{L_x}$ 个态、每格点 $2$ 个。故边缘激发由一个马约拉纳费米子描写：$L_x$ 个马约拉纳算符的不可约表示维数恰为 $2^{L_x/2}$。

\section*{问题 10.3.1}
\begin{problem}
（a）证明各 $\hat F_{\pmb i}$（$\hat F_{\pmb i}=\sigma_{\pmb i}^{x}\sigma_{\pmb i+\hat{\pmb x}}^{y}\sigma_{\pmb i+\hat{\pmb x}+\hat{\pmb y}}^{x}\sigma_{\pmb i+\hat{\pmb y}}^{y}$，式 (10.3.2)）在两方向均为周期性边界条件的格子上满足 $\prod_{\pmb i}\hat F_{\pmb i}=1$（式 (10.3.3)）。（b）在偶 $\times$ 偶周期格子上证明更强的条件 $\prod_{\pmb i=\mathrm{even}}\hat F_{\pmb i}=1$（式 (10.3.4)）。
\end{problem}
\solution
\textbf{（a）}每个格点 $\pmb i$ 是四个方格的角：$\pmb i$（左下角，贡献 $\sigma_{\pmb i}^{x}$）、$\pmb i-\hat{\pmb x}$ 的右下角 $\pmb i+\hat{\pmb x}$（贡献 $\sigma_{\pmb i}^{y}$）、$\pmb i-\hat{\pmb y}$ 的左上角 $\pmb i+\hat{\pmb y}$（贡献 $\sigma_{\pmb i}^{y}$）、$\pmb i-\hat{\pmb x}-\hat{\pmb y}$ 的右上角 $\pmb i+\hat{\pmb x}+\hat{\pmb y}$（贡献 $\sigma_{\pmb i}^{x}$）。把乘积 $\prod_{\pmb p}\hat F_{\pmb p}$ 按任意固定次序排列，并将同一格点的四个泡利算符聚拢（不同格点的泡利算符对易，移位无代价）：每个格点上得到 $\sigma_{\pmb i}^{x}\sigma_{\pmb i}^{y}\sigma_{\pmb i}^{y}\sigma_{\pmb i}^{x}$（或其一重新排列）。聚拢过程中的同格点换序数：四个算符为两个 $\sigma^{x}$、两个 $\sigma^{y}$，把它们排成 $(x,x,y,y)$ 所需的同格点对换个数为偶数（$0$ 或 $2$，取决于到达次序），符号 $+1$。于是
\begin{equation*}
\prod_{\pmb p}\hat F_{\pmb p}=\prod_{\pmb i}\big(\sigma_{\pmb i}^{x}\big)^{2}\big(\sigma_{\pmb i}^{y}\big)^{2}=1 .
\end{equation*}
此论证只用到周期性边界条件保证每格点恰为四个方格的角，对格子的奇偶没有任何要求：
\begin{equation*}
\boxed{\;\prod_{\pmb i}\hat F_{\pmb i}=1\;}
\end{equation*}
即式 (10.3.3)。（这使通量组态 $\{F_{\pmb i}\}$ 不是独立的：$F$ 的每一种可取组态都必须满足 $\prod_{\pmb i}F_{\pmb i}=1$。）

\textbf{（b）}设 $L_x,L_y$ 均为偶数，只对偶方格（$(x_p+y_p)$ 偶）求积。此时每个格点是哪些偶方格的角？
\begin{itemize}
\item 偶格点 $\pmb i$（$x_i+y_i$ 偶）：方格 $\pmb i$（偶，左下角，$\sigma^{x}$）与 $\pmb i-\hat{\pmb x}-\hat{\pmb y}$（偶，右上角，$\sigma^{x}$）；而 $\pmb i-\hat{\pmb x}$、$\pmb i-\hat{\pmb y}$ 均为奇方格，不参与。故偶格点恰贡献两个 $\sigma^{x}$。
\item 奇格点 $\pmb i$：作为偶方格 $\pmb i-\hat{\pmb x}$ 的右下角与 $\pmb i-\hat{\pmb y}$ 的左上角（两处均 $\sigma^{y}$），恰贡献两个 $\sigma^{y}$。
\end{itemize}
（校验计数：偶方格 $N_{\text{site}}/2$ 个 $\times$ 四角 $=2N_{\text{site}}$ 个角位；偶格点 $N_{\text{site}}/2\times2$ 加奇格点 $N_{\text{site}}/2\times2=2N_{\text{site}}$，自洽。）于是
\begin{equation*}
\prod_{\pmb p=\mathrm{even}}\hat F_{\pmb p}=\prod_{\pmb i\ \text{偶}}\big(\sigma_{\pmb i}^{x}\big)^{2}\prod_{\pmb i\ \text{奇}}\big(\sigma_{\pmb i}^{y}\big)^{2}=1,
\end{equation*}
即式 (10.3.4)：
\begin{equation*}
\boxed{\;\prod_{\pmb i=\mathrm{even}}\hat F_{\pmb i}=1\qquad(\text{偶}\times\text{偶格子}),\qquad\text{同理}\quad\prod_{\pmb i=\mathrm{odd}}\hat F_{\pmb i}=1\;}
\end{equation*}
（第二式由 (a) 除以第一式得到）。这两条约束使通量组态数降为 $2^{N_{\text{site}}-2}$，从而如正文 10.3.2 节所述，$\hat F_{\pmb i}$ 的共同本征态必为四重简并。若格子不是偶 $\times$ 偶（例如 $4\times3$），上述角位配对失效，$\prod_{\mathrm{even}}\hat F_{\pmb i}=\prod_{\text{偶格点}}\sigma_{\pmb i}^{x}\sigma_{\pmb i}^{y}$ 不再是 $\pm1$ 型恒等式——这正是 (b) 限于偶 $\times$ 偶格子的原因。

\section*{问题 10.5.1}
\begin{problem}
（a）证明式 (10.5.3) 中的 $4\times4$ 矩阵 $\gamma^{a}$ 满足狄拉克代数 $\{\gamma^{a},\gamma^{b}\}=2\delta_{ab}$。（b）证明式 (10.5.4) 求和中的各项彼此对易（提示：可用 $\gamma^{ab}$ 的马约拉纳费米子表示），从而使 $H_{3D}$ 可精确求解。
\end{problem}
\solution
\textbf{（a）}$\gamma^{a}\equiv\gamma^{az}$（$a=x,\bar x,y,\bar y$）的四个矩阵为
\begin{equation*}
\gamma^{x}=\sigma^{x}\otimes\sigma^{x},\quad
\gamma^{\bar x}=\sigma^{y}\otimes\sigma^{x},\quad
\gamma^{y}=\sigma^{z}\otimes\sigma^{x},\quad
\gamma^{\bar y}=\sigma^{0}\otimes\sigma^{y}.
\end{equation*}
对角项：$(\gamma^{a})^{2}=1$（$\sigma^{a},\sigma^{0}$ 均平方为 $1$）。对非对角项逐一检验，前因子或后因子必有一个反对易：
\begin{align*}
\gamma^{x}\gamma^{\bar x}&=(\sigma^{x}\sigma^{y})\otimes(\sigma^{x}\sigma^{x})=\mathrm{i}\sigma^{z}\otimes 1,&
\gamma^{\bar x}\gamma^{x}&=-\,\mathrm{i}\sigma^{z}\otimes1,\\
\gamma^{x}\gamma^{y}&=(\sigma^{x}\sigma^{z})\otimes(\sigma^{x}\sigma^{x})=-\mathrm{i}\sigma^{y}\otimes1,&
\gamma^{y}\gamma^{x}&=+\mathrm{i}\sigma^{y}\otimes1,\\
\gamma^{x}\gamma^{\bar y}&=(\sigma^{x}\sigma^{0})\otimes(\sigma^{x}\sigma^{y})=\mathrm{i}\,\sigma^{x}\otimes\sigma^{z},&
\gamma^{\bar y}\gamma^{x}&=-\mathrm{i}\,\sigma^{x}\otimes\sigma^{z},\\
\gamma^{y}\gamma^{\bar y}&=(\sigma^{z}\sigma^{0})\otimes(\sigma^{x}\sigma^{y})=-\mathrm{i}\,\sigma^{z}\otimes\sigma^{z},&
\gamma^{\bar y}\gamma^{y}&=+\mathrm{i}\,\sigma^{z}\otimes\sigma^{z},\\
\gamma^{\bar x}\gamma^{y}&=(\sigma^{y}\sigma^{z})\otimes(\sigma^{x}\sigma^{x})=\mathrm{i}\sigma^{x}\otimes1,&
\gamma^{y}\gamma^{\bar x}&=-\mathrm{i}\sigma^{x}\otimes1,\\
\gamma^{\bar x}\gamma^{\bar y}&=(\sigma^{y}\sigma^{0})\otimes(\sigma^{x}\sigma^{y})=\mathrm{i}\,\sigma^{y}\otimes\sigma^{z},&
\gamma^{\bar y}\gamma^{\bar x}&=-\mathrm{i}\,\sigma^{y}\otimes\sigma^{z}.
\end{align*}
每一对的和均为零，且对角项给 $2\delta_{ab}$，故 $\{\gamma^{a},\gamma^{b}\}=2\delta_{ab}$，即狄拉克代数。

\textbf{（b）}用马约拉纳表示 $\gamma^{ab}=\mathrm{i}\lambda^{a}\lambda^{b}$（$a\neq b$；同格点指标略去），式 (10.5.4) 的每个方格项是 $8$ 个马约拉纳算符的乘积——每个角上两个，分别来自该角的两个方向的链（链与指标对的对应见图 10.7：$+\hat x$ 链带 $(x,\bar x)$、$+\hat y$ 链带 $(y,\bar y)$、$+\hat z$ 链带 $(z,\bar z)$）。以 $xy$ 方格为例，
\begin{equation*}
-\gamma_{\pmb i}^{yx}\gamma_{\pmb{i}+\hat{\pmb x}}^{\bar x y}\gamma_{\pmb{i}+\hat{\pmb x}+\hat{\pmb y}}^{\bar y\bar x}\gamma_{\pmb{i}+\hat{\pmb y}}^{x\bar y}
\;\sim\;
\big(\lambda_{\pmb i}^{x}\lambda_{\pmb i}^{y}\big)\,
\big(\lambda_{\pmb{i}+\hat{\pmb y}}^{x}\lambda_{\pmb{i}+\hat{\pmb y}}^{\bar y}\big)\,
\big(\lambda_{\pmb{i}+\hat{\pmb x}}^{\bar x}\lambda_{\pmb{i}+\hat{\pmb x}}^{y}\big)\,
\big(\lambda_{\pmb{i}+\hat{\pmb x}+\hat{\pmb y}}^{\bar x}\lambda_{\pmb{i}+\hat{\pmb x}+\hat{\pmb y}}^{\bar y}\big),
\end{equation*}
即四个角双线性因子之积（乘积中每个角恰出现一次，个别方的最远角可出现两次而相互抵消，不影响下述计数）。任意两个方格项 $\hat F_{\pmb p},\hat F_{\pmb p'}$ 的对易性归结为角双线性因子的换序：同一格点上两个双线性因子 $\lambda^{a}\lambda^{b}$ 与 $\lambda^{c}\lambda^{d}$ 换序时，若共享一个马约拉纳算符则给出 $(-1)$，否则给出 $+1$（$2\times2=4$ 次同格点对换），若相同则平方为 $+1$。按两方格的相对位置分类：
\begin{itemize}
\item 不相交：平凡对易；
\item 共享一条边（两个共享格点，各自恰有一个共享味）：每个共享格点贡献一个 $(-1)$，共 $(-1)^{2}=+1$；
\item 只共享一个角格点：该格点上四个马约拉纳算符分属四条不同的链（立方格子每格点六条链的指标对互不相同），换序符号 $+1$。
\end{itemize}
故所有 $[\hat F_{\pmb p},\hat F_{\pmb p'}]=0$。同理 $\hat F_{\pmb p}^{2}=1$（每方格四个角双线性各自平方为 $-1$，共 $(-1)^{4}=+1$），且每个 $\hat F_{\pmb p}$ 与所有链算符对易。于是 $H_{3D}=-g\sum_{\pmb p}\hat F_{\pmb p}$ 由共同本征态 $|\{F_{\pmb p}\},\alpha\rangle$ 对角化，能量 $-g\sum_{\pmb p}F_{\pmb p}$，精确可解。
\emph{（译注：原书式 (10.5.4) 第三项的最后一个因子印作 $\gamma_{\pmb{i}+\hat{\pmb x}}^{z\bar y}$；按指标对约定（$+\hat x$ 链带 $(x,\bar x)$、$+\hat z$ 链带 $(z,\bar z)$），$xz$ 方格第四个角上的指标对应为 $(z,\bar x)$，即应为 $\gamma_{\pmb{i}+\hat{\pmb x}}^{z\bar x}$，$\bar y$ 为 $\bar x$ 之误排。上文按此改正。）}

\section*{问题 10.5.2}
\begin{problem}
用马约拉纳费米子与投影构造求解式 (10.5.4)：(a) 证明式 (10.5.5) 定义的诸 $\hat U_{\pmb{i}\pmb j}$ 相互对易；(b) 证明在每格点费米子数为偶的子空间内，费米子哈密顿量 (10.5.6) 约化为自旋 $3/2$ 哈密顿量 (10.5.4)；(c) 用 10.2 节的程序求偶 $\times$ 偶 $\times$ 偶周期格子上的基态能量与基态简并度。
\end{problem}
\solution
\textbf{（a）}$\hat U_{\pmb{i}\pmb j}=\lambda_{\pmb i}^{a}\lambda_{\pmb j}^{b}$，其中 $ab$ 是链 $\pmb i\to\pmb j$ 的指标对（$+\hat x$ 链带 $(x,\bar x)$，等等；见图 10.7）。与二维情形（问题 10.2.1）完全相同的计算给出：
\begin{enumerate}
\item[(i)] $\hat U_{\pmb{i}\pmb j}^{2}=\lambda_{\pmb i}^{a}\lambda_{\pmb j}^{b}\lambda_{\pmb i}^{a}\lambda_{\pmb j}^{b}=+1$（不同格点对易；注意此处无 $\mathrm i$ 因子，见下面译注）；
\item[(ii)] 不共享格点的两条链算符对易；
\item[(iii)] 共享一个格点的两条链算符反对易（两链在该格点的指标对不同）。
\end{enumerate}
\emph{（译注：与问题 10.2.1(a) 相同，"诸 $\hat U_{\pmb{i}\pmb j}$ 构成对易集合"的表述欠准确——相邻链算符实际上反对易。可解性所需且成立的是：方格算符 $\hat F_{\pmb p}$ 彼此对易（见问题 10.5.1(b)），每个 $\hat F_{\pmb p}$ 与每个 $\hat U_{\pmb{i}\pmb j}$ 对易（链 $\langle\pmb i\pmb j\rangle$ 与含它的两个方格共享该链的两个马约拉纳算符、与其它方格至多共享两个格点，换序符号均为偶），因而 $H=-\sum_{\pmb p}g\hat F_{\pmb p}$ 与所有 $\hat U_{\pmb{i}\pmb j}$ 对易。）}

\textbf{（b）}复费米子 $2\psi_{a}=\lambda^{a}+\mathrm{i}\lambda^{\bar a}$（$a=x,y,z$）满足标准反对易关系（同问题 10.2.1(b)），每格点 $8$ 维费米子空间；偶（奇）费米子数的子空间各为 $4$ 维，恰与自旋 $3/2$ 的 $4$ 维空间对应。把 $\gamma^{ab}=\mathrm{i}\lambda^{a}\lambda^{b}$ 代入 (10.5.4)：每个方格项是四个 $\gamma$ 之积，即 $8$ 个马约拉纳算符的乘积（$\mathrm i^{4}=1$），其马约拉纳内容恰为围绕方格四条链的四个 $\hat U_{\pmb i\pmb j}$ 之积的规范排列（逐链核对指标对，见图 10.7）。计入取向约定 $\hat U_{\pmb j\pmb i}=-\hat U_{\pmb i\pmb j}$ 与每个 $\gamma^{ab}$ 中的相因子 $\mathrm i$（$i^{4}=1$）带来的净符号 $(-1)$，得到作为算符恒等式的
\begin{equation*}
\hat F_{\pmb p}=\hat U_{\pmb i_1\pmb i_2}\hat U_{\pmb i_2\pmb i_3}\hat U_{\pmb i_3\pmb i_4}\hat U_{\pmb i_4\pmb i_1}
=-\,\gamma^{a_1b_1}\gamma^{a_2b_2}\gamma^{a_3b_3}\gamma^{a_4b_4}
\end{equation*}
（以 $xy$ 方格为例：$-\gamma^{yx}\gamma^{\bar x y}\gamma^{\bar y\bar x}\gamma^{x\bar y}$，与式 (10.5.4) 中 $\hat F_{\pmb p}=-\gamma^{yx}\gamma^{\bar x y}\gamma^{\bar y\bar x}\gamma^{x\bar y}$ 的定义一致）。因此费米子哈密顿量 (10.5.6) 与自旋 $3/2$ 哈密顿量 (10.5.4) 在全希尔伯特空间上恒同，特别地在每格点偶费米子的子空间内，它就是立方格子上的自旋 $3/2$ 模型（$\gamma^{ab}$ 正是作用在自旋 $3/2$ 四重态上的算符）。与二维情形不同的是：此处不投影即已恒等，投影只用于说明 $\hat F_{\pmb p}$ 在偶、奇两个子空间内给出两份全同的自旋 $3/2$ 系统。

\textbf{（c）}由 $[\hat F_{\pmb p},\hat F_{\pmb p'}]=0$、$\hat F_{\pmb p}^{2}=1$，本征态为 $|\{F_{\pmb p}\},\alpha\rangle$，能量 $E=-g\sum_{\pmb p}F_{\pmb p}$。周期格子有 $3N_{\text{site}}$ 个方格，故
\begin{equation*}
\boxed{\;E_{0}=-3\,|g|\,N_{\text{site}}\;}
\end{equation*}
（基态取 $F_{\pmb p}=\mathrm{sgn}(g)$ 对所有 $\pmb p$；每个激发为翻转若干 $F_{\pmb p}$，能量 $E_0+2|g|\times$（翻转数），且受通量守恒约束：每个单胞立方体满足 $\prod_{\pmb p\in S}F_{\pmb p}=1$）。基态简并度用 10.2 节的程序计数：物理态由链变量组态 $s_{\pmb{i}\pmb j}=\pm\mathrm i$（每条链一个 $Z_2$ 变量，共 $3N_{\text{site}}$ 个）的 $Z_2$ 规范等价类标记（$s_{\pmb{i}\pmb j}\to G_{\pmb i}s_{\pmb{i}\pmb j}G_{\pmb j}$，有效规范群 $2^{N_{\text{site}}-1}$），并带总费米子奇偶约束（三维版的 (10.2.7)--(10.2.8)：全部链算符之积给出各格点奇偶算符之积）。基态要求全部 $3N_{\text{site}}$ 个方格通量 $F_{\pmb p}=\mathrm{sgn}(g)$。通量映射（链变量 $\to$ 面通量）的秩为 $2N_{\text{site}}-2$（其核由 $N_{\text{site}}-1$ 个立方体约束与三个全局取向约束张成），故固定通量组态的组态数为 $2^{3N_{\text{site}}-(2N_{\text{site}}-2)}=2^{N_{\text{site}}+2}$；除以规范等价 $2^{N_{\text{site}}-1}$，得
\begin{equation*}
\boxed{\;\text{基态简并度}=2^{N_{\text{site}}+2-(N_{\text{site}}-1)}=2^{3}=8\;}
\end{equation*}
这三个二重选择正是三维环面三个不可收缩方向上各自的 $Z_2$ 和乐（holonomy）：与二维环面上的四重简并 $2^2$ 完全平行，且与立方格子上 $Z_2$ 规范理论在 $T^3$ 上的基态简并度 $|H^1(T^3;Z_2)|=2^3=8$ 一致。这些基态间的跃迁算符是非局域的（局域地看与 $Z_2$ 规范变换不可区分），因而简并对任意局域微扰稳健。

\section*{问题 10.5.3}
\begin{problem}
（a）利用 $\gamma^{ab}=\mathrm{i}\lambda^{a}\lambda^{b}$，证明闭合弦算符 (10.5.7) 可表示为 $W(C_{\text{close}})=\hat U_{\pmb i_1\pmb i_2}\hat U_{\pmb i_2\pmb i_3}\cdots\hat U_{\pmb i_{n-1}\pmb i_n}\hat U_{\pmb i_n\pmb i_1}$。（b）证明任意两个闭合弦算符彼此对易。（c）验证关系式 (10.5.8)：$W(C)=W(C_1)W(C_2)$。
\end{problem}
\solution
\textbf{（a）}关键为 $-\mathrm{i}\gamma_{\pmb i}^{ab}=\lambda_{\pmb i}^{a}\lambda_{\pmb i}^{b}$。闭合弦 $C$ 依次穿过格点 $\pmb i_1,\dots,\pmb i_n$，且 $b_ma_{m+1}$ 是链 $\pmb i_m\to\pmb i_{m+1}$ 的指标对。由 (10.5.7)，
\begin{equation*}
W(C)=-\prod_{m=1}^{n}\big(-\mathrm{i}\gamma_{\pmb i_m}^{a_mb_m}\big)
=-\prod_{m=1}^{n}\lambda_{\pmb i_m}^{a_m}\lambda_{\pmb i_m}^{b_m},
\end{equation*}
（由 $\gamma^{ab}=\mathrm{i}\lambda^{a}\lambda^{b}$ 得 $-\mathrm i\gamma^{ab}=\lambda^{a}\lambda^{b}$；开头的整体 $-1$ 即 (10.5.7) 前的负号）。另一方面，按 (10.5.5)，$\hat U_{\pmb i_m\pmb i_{m+1}}=\lambda_{\pmb i_m}^{b_m}\lambda_{\pmb i_{m+1}}^{a_{m+1}}$，于是
\begin{equation*}
\prod_{m=1}^{n}\hat U_{\pmb i_m\pmb i_{m+1}}
=\big(\lambda_{\pmb i_1}^{b_1}\lambda_{\pmb i_2}^{a_2}\big)\big(\lambda_{\pmb i_2}^{b_2}\lambda_{\pmb i_3}^{a_3}\big)\cdots\big(\lambda_{\pmb i_n}^{b_n}\lambda_{\pmb i_1}^{a_1}\big),
\end{equation*}
其马约拉纳多重集与 $W(C)$ 相同（每个弦格点上恰有 $\lambda^{a_m},\lambda^{b_m}$ 各一次）。把末端的 $\lambda_{\pmb i_1}^{a_1}$ 移到最前：它跨过 $2n-1$ 个算符，其中唯一的同格点配对是 $\lambda_{\pmb i_1}^{b_1}$，给出唯一的 $(-1)$，其余移动均在不同格点间进行（无符号）。故
\begin{equation*}
\prod_{m}\hat U_{\pmb i_m\pmb i_{m+1}}=-\prod_{m}\lambda_{\pmb i_m}^{a_m}\lambda_{\pmb i_m}^{b_m}=W(C_{\text{close}}),
\end{equation*}
即
\begin{equation*}
\boxed{\;W(C_{\text{close}})=\hat U_{\pmb i_1\pmb i_2}\hat U_{\pmb i_2\pmb i_3}\cdots\hat U_{\pmb i_{n-1}\pmb i_n}\hat U_{\pmb i_n\pmb i_1}\;}
\end{equation*}

\textbf{（b）}由（a），两个闭合弦算符都是马约拉纳双线性因子的乘积（每个弦格点一个双线性因子）。同一格点上两个双线性因子 $\lambda^{a}\lambda^{b}$ 与 $\lambda^{c}\lambda^{d}$ 的换序规则：共享一个马约拉纳算符时反号，不共享时 $2\times2=4$ 次对换无符号，相同时平方无符号。考察两个闭合弦 $C,C'$ 的关系：
\begin{itemize}
\item 若它们不相交或只交于一个格点（不共享链）：相交格点上两个双线性因子的四个马约拉纳算符来自四条不同链，互不共享，换序符号 $+1$；
\item 若它们共享一段路径（若干条公共链）：路径的中间格点上两个双线性因子完全相同，对易是平凡的；而路径的两个端点上两个双线性因子恰共享一个马约拉纳算符（公共链的味），各贡献一个 $(-1)$，总共 $(-1)^{2}=+1$。
\end{itemize}
故任意两个闭合弦算符对易。（这也解释了正文 10.5.2 节的论断：闭合弦算符与 $\hat F_{\pmb p}$——它本身就是一个闭合弦算符——对易，从而基态具有闭合弦凝聚。）

\textbf{（c）}设回路 $C$ 分解为 $C_1\cup C_2$，$C_1$ 与 $C_2$ 共享一段路径 $P$，且 $C$ 是 $C_1$ 与 $C_2$ 的对称差。由（a），三个闭合弦算符都是沿各自回路的链算符 $\hat U$ 之积。在乘积 $W(C_1)W(C_2)$ 中，公共路径 $P$ 上的每个链算符 $\hat U_{\pmb{i}\pmb j}$ 出现两次；把两个乘积按回路序连接后，成对出现的相同马约拉纳算符相邻化并湮灭（$\lambda^{2}=1$，成对移动的换序符号成对抵消），剩下的正是沿 $C$ 的链算符之积。以正方形并排的 $1\times2$ 矩形回路为例（$C$ 为外边界，$C_1,C_2$ 为两个方格边界，共享一条内部边）：$C$ 有 $6$ 条边、$C_1,C_2$ 各 $4$ 条边；$W(C_1)W(C_2)$ 中公共边的两个 $\hat U$ 相乘为零（$\hat U^2$ 贡献常数 $+1$，两条公共链共 $(-1)^2=+1$），剩余链恰好沿 $C$ 排列，逐项比较给出
\begin{equation*}
\boxed{\;W(C)=W(C_1)W(C_2)\;}
\end{equation*}
即式 (10.5.8)。这一关系使任何闭合弦算符可写成 $\hat F_{\pmb p}$ 的乘积（把回路取为方格面的组合边界），从而可计算凝聚振幅：对 $F_{\pmb p}=+1$ 的基态 $\langle W(C_{\text{close}})\rangle=1$；对 $F_{\pmb p}=-1$ 的基态 $\langle W(C_{\text{close}})\rangle=(-1)^{N_{\pmb p}}$，$N_{\pmb p}$ 为以 $C$ 为边界的曲面上的方格数。

\section*{问题 10.6.1}
\begin{problem}
考虑式 (10.6.3) 中 $i=1,2,3$ 的三转子模型（图 10.11）：三个转子放在三角形的链 $\langle12\rangle$、$\langle23\rangle$、$\langle31\rangle$ 上，
\begin{equation*}
H=\sum_{i=1}^{3}\Big[U\big(S_{\langle i-1,i\rangle}^{z}-S_{\langle i,i+1\rangle}^{z}\big)^{2}+J\big(S_{\langle i,i+1\rangle}^{z}\big)^{2}+t\big(\mathrm e^{\mathrm i\theta_{\langle i-1,i\rangle}}\mathrm e^{\mathrm i\theta_{\langle i,i+1\rangle}}+\text{h.c.}\big)\Big].
\end{equation*}
（1）导出低能格点规范理论的作用量（$g$ 只需到 $O(1)$ 系数精度）；（2）设 $U\gg g\gg J$，计算低能能级。
\end{problem}
\solution
\textbf{（1）低能规范理论的导出。}$\theta_{\langle i,i+1\rangle}$ 与 $S_{\langle i,i+1\rangle}^{z}$ 是正则共轭对，路径积分为
\begin{equation*}
Z=\int\mathcal D S^{z}\,\mathcal D\theta\;
\mathrm e^{\mathrm i\int\mathrm dt\,\big[\sum_{i}S_{\langle i,i+1\rangle}^{z}\dot\theta_{\langle i,i+1\rangle}-H\big]} .
\end{equation*}
引入 $a_0(i)$（$i=1,2,3$）解耦 $U$ 项：$ab=\dfrac{(a+b)^2}{4U}-\dfrac{(a-b)^2}{4U}$ 型恒等式给出
\begin{equation*}
U\big(S_{i-1}^{z}-S_{i}^{z}\big)^{2}\;\longrightarrow\;a_0(i)\big(S_{i-1}^{z}-S_{i}^{z}\big)-\frac{a_0^{2}(i)}{4U},
\end{equation*}
对 $S_{\langle i,i+1\rangle}^{z}$（高斯型积分）完成积分后得到拉格朗日量
\begin{equation*}
L=\frac{1}{4J}\sum_{i=1}^{3}\Big[\big(\dot\theta_{\langle i,i+1\rangle}+a_0(i)-a_0(i+1)\big)^{2}+\frac{a_0^{2}(i)}{4U}\Big]
-t\sum_{i}\big(\mathrm e^{\mathrm i(\theta_{\langle i-1,i\rangle}+\theta_{\langle i,i+1\rangle})}+\text{h.c.}\big).
\end{equation*}
在大 $U$ 极限下舍去 $a_0^{2}/4U$ 项，识别 $a_{i,i+1}\equiv\theta_{\langle i,i+1\rangle}$（$a_{i+1,i}=-a_{i,i+1}$）为链上的格点规范场：
\begin{equation*}
L=\frac{1}{4J}\sum_{i=1}^{3}\big(\dot a_{i,i+1}+a_0(i)-a_0(i+1)\big)^{2}
-t\sum_{i}\big(\mathrm e^{\mathrm i(a_{i-1,i}+a_{i,i+1})}+\text{h.c.}\big)+\frac{a_0^{2}}{4U},
\end{equation*}
它在规范变换 $a_{ij}\to a_{ij}+\phi_j-\phi_i$、$a_0(i)\to a_0(i)+\dot\phi_i$ 下不变。$U$ 项强制低能态满足 $S_{12}=S_{23}=S_{31}=n$（能隙 $O(U)$），低能态为 $|nnn\rangle$，$E_n=3Jn^{2}$；低能波函数 $\Psi(a_{12},a_{23},a_{31})$ 是 $|nnn\rangle$ 的叠加且规范不变。

$t$ 项的处理与四转子情形相同：纯规范涨落 $a_{ij}=\phi_j-\phi_i$ 使 $\mathrm e^{\mathrm i(a_{i-1,i}+a_{i,i+1})}=\mathrm e^{\mathrm i(\phi_{i+1}-\phi_{i-1})}$ 以 $O(U)$ 的能隙快涨落，$t$ 项一阶平均为零。因此规范不变的势能项只能由多阶过程产生：规范不变的字面组合须是 $\mathrm e^{\pm\mathrm i\Phi}$，$\Phi=a_{12}+a_{23}+a_{31}$ 为穿过三角形的总通量。这里出现一个与四转子（四方格）本质不同的代数约束：$t$ 顶点 $\mathrm e^{\mathrm i(a_{i-1,i}+a_{i,i+1})}$ 使 $\sum_i S_{\langle i,i+1\rangle}^{z}$ 改变 $\pm2$，故
\begin{equation*}
\big[(-1)^{S_{12}^{z}+S_{23}^{z}+S_{31}^{z}},\,H\big]=0
\end{equation*}
是精确的守恒律；而 $\mathrm e^{\mathrm i\Phi}$ 使 $\sum S^{z}$ 改变 $-3$（奇数），$|nnn\rangle\to|n_1n_1n_1\rangle$（$n_1=n+1$）在任何有限阶微扰中都被禁戒。事实上二阶过程 $|nnn\rangle\to|n,n\pm1,n\pm1\rangle\to\cdots$ 只能回到 $|nnn\rangle$ 自身（六个中间态各贡献 $-t^{2}/2U$），故
\begin{equation*}
\delta E_n=-\frac{6t^{2}}{2U}=-\frac{3t^{2}}{U},
\end{equation*}
而 $\langle n_1n_1n_1|H_{\rm eff}|nnn\rangle=0$ 严格成立：三角形是奇数环，$O(t^{2}/U)$ 的 $g\cos\Phi$ 项不存在。首个规范不变的非平庸项出现在第三阶：三个 $t$ 顶点之积 $t^{3}\prod_{i}\mathrm e^{\mathrm i(a_{i-1,i}+a_{i,i+1})}=t^{3}\mathrm e^{2\mathrm i\Phi}$ 与 $\phi_i$ 无关、不经受平均，故低能有效拉格朗日量为
\begin{equation*}
\boxed{\;L=\frac{1}{4J}\sum_{i=1}^{3}\big(\dot a_{i,i+1}+a_0(i)-a_0(i+1)\big)^{2}+g_{3}\cos(2\Phi),
\qquad g_{3}=O\!\Big(\frac{t^{3}}{U^{2}}\Big),\quad \Phi=\sum_i a_{i,i+1}\;}
\end{equation*}
即仍是一个格点 $U(1)$ 规范理论，但有效势的周期为 $\pi$（$\Phi$ 的 $Z_2$ 剩余对称性是奇数环上 $\sum S^{z}$ 宇称守恒的场论体现）。
\emph{（译注：若不加分辨地照搬四方格的结果，会写出 $g=O(t^{2}/U)$ 的 $g\cos\Phi$；由上述精确宇称选择定则，该项在三角形上恒为零，$g$ 的领头项是 $O(t^{3}/U^{2})$ 的 $\cos2\Phi$。若只问"是否有低能规范理论"，答案不变：动能项加规范不变的势能项，规范群 $U(1)$。）}

\textbf{（2）低能能级（$U\gg g\gg J$）。}令 $U\to\infty$，取规范 $a_0=0$，用均匀模 $\Phi=\sum_i a_{i,i+1}$（三个纯规范组合 $\phi_i$ 有 $O(U)$ 能隙，可积掉）。动能项：
\begin{equation*}
\frac{1}{4J}\sum_{i}\dot a_{i,i+1}^{2}\;\supset\;\frac{\dot\Phi^{2}}{12J}
\qquad\Longrightarrow\qquad
H_{\Phi}=3J\,P_{\Phi}^{2}+\big(\text{势能}\big),
\end{equation*}
（$P_{\Phi}=\partial L/\partial\dot\Phi=\dot\Phi/6J$；检验：$g=0$ 时能级 $E_m=3Jm^{2}$，与 $|nnn\rangle$ 的能量 $3Jn^{2}$ 精确一致，这正是正文对四转子用过的自洽检验的三角形版本）。以领头势能 $V(\Phi)=g\cos(2\Phi)$（$g\equiv g_3$）为例：$g\gg J$ 为半经典极限，低能级是 $\Phi=0,\pi$ 两个势阱附近的摆动态：
\begin{equation*}
V(\Phi)\approx-g+2g\,(2\,\delta\Phi)^{2}/2=-g+4g\,\delta\Phi^{2},
\end{equation*}
谐振子频率 $\omega=2\sqrt{3J\cdot4g}=4\sqrt{3Jg}$，故
\begin{equation*}
\boxed{\;E_{m}\approx-g+4\sqrt{3Jg}\,\Big(m+\tfrac12\Big)+O\!\Big(\frac{J^{2}}{g^{1/2}}\Big),\qquad m=0,1,2,\dots\;}
\end{equation*}
两个势阱（$\Phi=0,\pi$）由 $\pm\Phi$ 的绕数跃迁（即 $\mathrm e^{\pm\mathrm i\Phi}$ 型过程，被宇称禁戒、以指数小的高阶效应出现）沟通，故低能级在此精度下二重简并；更高的能级为绕数态，$E\sim3J(m')^{2}$。若按四转子类比取 $g\cos\Phi$，同样的展开给出 $E_m\approx-g+\sqrt{6Jg}\,(m+\tfrac12)$，结构相同。

\section*{问题 10.6.2}
\begin{problem}
验证式 (10.6.16) 与 (10.6.17)：二维与三维转子格子模型在大 $U$ 极限下的低能有效拉格朗日量 (10.6.9)/(10.6.11) 的连续极限分别为
\begin{equation*}
L=\int\mathrm d^{2}\pmb x\;\Big(\frac{\pmb e^{2}}{4J}-\frac{gl^{2}}{2}\pmb b^{2}\Big),\qquad
L=\int\mathrm d^{3}\pmb x\;\Big(\frac{\pmb e^{2}}{4Jl}-\frac{gl}{2}\pmb b^{2}\Big),
\end{equation*}
其中 $\pmb e=\partial_t\pmb a-\nabla a_0$，$\pmb b=\nabla\times\pmb a$，$l$ 为晶格常数。
\end{problem}
\solution
\textbf{二维式 (10.6.16)。}出发点是式 (10.6.9)：
\begin{equation*}
L=\frac{1}{4J}\sum_{\pmb i,\pmb\mu=\pmb x,\pmb y}\big[\dot a_{\pmb i,\pmb i+\pmb\mu}+a_0(\pmb i)-a_0(\pmb i+\pmb\mu)\big]^{2}
+g\sum_{\pmb i}\cos\Phi_{\pmb i},
\end{equation*}
其中 $a_{\pmb i\pmb j}=-a_{\pmb j\pmb i}$、$a_{\pmb i,\pmb i+\pmb\mu}=(-)^{\pmb i}\theta_{\pmb i+\pmb\mu/2}$。按连续极限
\begin{equation*}
a_{\pmb i,\pmb i+\pmb\mu}=\delta\pmb x_{\langle\pmb i\pmb j\rangle}\cdot\pmb a(\pmb x)=l\,a_\mu(\pmb x),\qquad
a_0(\pmb i)=a_0(\pmb x),
\end{equation*}
其中 $\pmb x$ 取在链的中点附近，$\delta\pmb x_{\langle\pmb i\pmb j\rangle}=\pmb\mu\,l$。

\emph{动能项。}由 $a_0(\pmb i)-a_0(\pmb i+\pmb\mu)=-l\,\partial_\mu a_0+O(l^{2})$ 及 $\dot a_{\pmb i,\pmb i+\pmb\mu}=l\,\partial_t a_\mu+O(l^{3})$：
\begin{equation*}
\big[\dot a_{\pmb i,\pmb i+\pmb\mu}+a_0(\pmb i)-a_0(\pmb i+\pmb\mu)\big]^{2}
=l^{2}\big(\partial_t a_\mu-\partial_\mu a_0\big)^{2}+O(l^{3}) .
\end{equation*}
求和 $\sum_{\pmb i,\pmb\mu}$ 中每个方向给出 $\sum_{\pmb i}l^{2}\to\int\mathrm d^{2}\pmb x$，故
\begin{equation*}
L_{\text{kin}}=\frac{1}{4J}\int\mathrm d^{2}\pmb x\;
\sum_{\mu}\big(\partial_t a_\mu-\partial_\mu a_0\big)^{2}
=\int\mathrm d^{2}\pmb x\;\frac{\pmb e^{2}}{4J},
\end{equation*}
其中 $\pmb e=\partial_t\pmb a-\nabla a_0$ 为仿造电场，与式 (10.6.16) 的第一项一致。

\emph{磁能项。}穿过方格 $\pmb i$ 的通量为
\begin{align*}
\Phi_{\pmb i}&=a_{\pmb i,\pmb i+\pmb x}+a_{\pmb i+\pmb x,\pmb i+\pmb x+\pmb y}+a_{\pmb i+\pmb x+\pmb y,\pmb i+\pmb y}+a_{\pmb i+\pmb y,\pmb i}\\
&=l\,a_x\big(\pmb x+\tfrac{\pmb x}{2}\big)+l\,a_y\big(\pmb x+\pmb x+\tfrac{\pmb y}{2}\big)
-l\,a_x\big(\pmb x+\pmb y+\tfrac{\pmb x}{2}\big)-l\,a_y\big(\pmb x+\tfrac{\pmb y}{2}\big)\\
&=l^{2}\big(\partial_x a_y-\partial_y a_x\big)+O(l^{3})=l^{2}\,b+O(l^{3}),
\end{align*}
即 $\Phi_{\pmb i}$ 是方格上的格点磁通量，连续极限下正比于仿造磁场 $b=\partial_xa_y-\partial_ya_x$。于是
\begin{equation*}
\sum_{\pmb i}g\cos\Phi_{\pmb i}
=\sum_{\pmb i}g\Big(1-\frac{\Phi_{\pmb i}^{2}}{2}+\cdots\Big)
\;\longrightarrow\;\text{常数}-\frac{1}{l^{2}}\int\mathrm d^{2}\pmb x\;\frac{g\,l^{4}}{2}\,b^{2},
\end{equation*}
（$\sum_{\pmb i}\to l^{-2}\int\mathrm d^{2}\pmb x$）。合并两项：
\begin{equation*}
\boxed{\;L=\int\mathrm d^{2}\pmb x\;\Big(\frac{\pmb e^{2}}{4J}-\frac{g\,l^{2}}{2}\,\pmb b^{2}\Big)\;}
\end{equation*}
即式 (10.6.16)：人工光的波速由 $\alpha_{\rm kin}=1/4J$、$\beta_{\rm mag}=gl^{2}/2$ 给出 $c_a=\sqrt{\beta/\alpha}=\sqrt{2gJ}\,l$（含 $\hbar$ 即 $c_a=\sqrt{2gJ}\,l^{2}/\hbar$），带宽 $E_a=2\hbar c_a/l=\sqrt{8gJ}$，与正文一致。

\textbf{三维式 (10.6.17)。}对式 (10.6.11)（求和遍及 $\pmb\mu=\pmb x,\pmb y,\pmb z$ 与全部方格取向）作同样的替换：动能项中 $\sum_{\pmb i,\pmb\mu}$ 的每个方向给出 $\sum_{\pmb i}l^{2}\to l^{-3}\int\mathrm d^{3}\pmb x\;l^{2}=l^{-1}\int\mathrm d^{3}\pmb x$：
\begin{equation*}
L_{\text{kin}}=\frac{1}{4J}\cdot\frac1l\int\mathrm d^{3}\pmb x\;\pmb e^{2}=\int\mathrm d^{3}\pmb x\;\frac{\pmb e^{2}}{4Jl}.
\end{equation*}
磁能项：$\Phi_{\pmb p}=l^{2}b_{\pmb p}+O(l^{3})$（$b_{\pmb p}$ 为该取向方格法向上的磁场分量），而 $\sum_{\pmb p}\to l^{-3}\int\mathrm d^{3}\pmb x$：
\begin{equation*}
\sum_{\pmb p}g_{\pmb p}\cos\Phi_{\pmb p}\;\longrightarrow\;\text{常数}-\frac{1}{l^{3}}\int\mathrm d^{3}\pmb x\;\frac{g\,l^{4}}{2}\,\pmb b^{2}
=-\int\mathrm d^{3}\pmb x\;\frac{g\,l}{2}\,\pmb b^{2}.
\end{equation*}
合并：
\begin{equation*}
\boxed{\;L=\int\mathrm d^{3}\pmb x\;\Big(\frac{\pmb e^{2}}{4Jl}-\frac{g\,l}{2}\,\pmb b^{2}\Big)\;}
\end{equation*}
即式 (10.6.17)。与式 (10.6.18) 的标准形式 $L=\int\mathrm d^{3}\pmb x\,(8\pi\alpha)^{-1}(\pmb e^{2}/c_a-c_a\pmb b^{2})$ 比较系数：
\begin{equation*}
\frac{1}{8\pi\alpha\,c_a}=\frac{1}{4Jl},\qquad
\frac{c_a}{8\pi\alpha}=\frac{gl}{2}
\;\Longrightarrow\;
c_a^{2}=2gJl^{2},\qquad
\alpha=\frac{1}{2\pi}\sqrt{\frac{J}{2g}},
\end{equation*}
与正文给出的 $c_a=\sqrt{2gJ}\,l^{2}/\hbar$ 和人工精细结构常数 $\alpha=(2\pi)^{-1}\sqrt{J/2g}$ 完全一致，验证完毕。（动能与磁能项中对 $l$ 的相反依赖 $1/l$ 与 $l$ 来自：电场能躺在链上（链数密度 $\sim l^{-3}$，每链贡献 $l^{2}$），而磁通量是面通量，$\Phi\sim l^{2}b$、方格密度 $\sim l^{-3}$。因此人工光的精细结构常数 $\alpha\sim\sqrt{J/g}$ 与晶格常数无关，而 $c_a\propto l$。）




